\chapter{Why There Is a Spread}\label{ch:mx:why-there-is-a-spread}

A market maker quotes 100.00 bid, 100.02 offered. If nobody knew more than she did, the two cents would pay for her time, her systems and the
risk of the inventory she carries between trades. A buyer who has read the news first makes the spread pay for something else: every trade
with him loses money, and the spread must recover the loss from everyone else. Three theories of the spread answer three questions: what it
costs to process a trade, what it costs to hold what a trade leaves behind, and what it costs to trade with someone who knows more. They lead
to different spreads, to different dynamics of quotes, and to different ways of measuring each part (chapter~5).

\section{Order-processing costs}

\begin{definition}[Order-processing cost]\label{def:mx:why-there-is-a-spread:processing}
The \emph{order-processing cost}\index{order-processing cost} of a liquidity provider is the cost per trade that does not depend on who the
counterparty is or on what happens to the price afterwards: exchange and clearing fees, systems, staff and capital, net of rebates.
\end{definition}

If processing were the only cost, a competitive market maker would quote a half-spread equal to it, the value would not react to trades, and
transaction prices would bounce between bid and ask around a random-walk value. Roll (1984) turned the bounce into an estimate: with a constant
effective half-spread $c$ and independent trade signs, the price changes have autocovariance $-c^2$, so $s = 2\sqrt{-\mathrm{Cov}(\Delta p_t,\Delta
p_{t-1})}$ (Roll's estimator, One Quant Book 4, chapter~21). On the simulated hour of \texttt{firm.tape}, the estimate is 0.88 ticks against a quoted
spread of 1.08: the simulator's trade signs are positively autocorrelated (0.29 at lag one, because noise orders come in metaorders), which removes
part of the bounce, as it does in real markets (One Quant Book 7, chapter~9).

\section{Inventory}

\begin{definition}[Inventory-holding cost]\label{def:mx:why-there-is-a-spread:inventory}
The \emph{inventory-holding cost}\index{inventory-holding cost} of a trade is the compensation a risk-averse liquidity provider requires for carrying
the position the trade leaves her with until she can unwind it.
\end{definition}

Garman (1976) modelled a dealer whose inventory is a random walk driven by the arrival of buyers and sellers, and who must set prices to avoid ruin;
Stoll (1978) and Ho and Stoll (1981) priced the risk. The simplest version already gives the two features every inventory model shares.

\begin{proposition}[Mean--variance dealer]\label{prop:mx:why-there-is-a-spread:inventory}
A dealer marks her inventory $q$ to the value $v$, expects to hold it for a time $T$ during which the value has variance $\sigma^2T$, and has mean--variance
preferences with risk aversion $\gamma$: holding $q$ costs her $\tfrac12\gamma\sigma^2Tq^2$. The prices at which a trade of size $Q$ leaves her indifferent
are
\[
a = v + \gamma\sigma^2T\left(\frac Q2 - q\right),\qquad b = v - \gamma\sigma^2T\left(\frac Q2 + q\right):
\]
a spread $\gamma\sigma^2TQ$ that grows with the size of the trade, centred on the reservation price $v-\gamma\sigma^2Tq$ that moves against her inventory.
\end{proposition}

\begin{proof}
Selling $Q$ at $a$ changes her wealth by $Q(a-v)$ and her holding cost from $\tfrac12\gamma\sigma^2Tq^2$ to $\tfrac12\gamma\sigma^2T(q-Q)^2$;
indifference gives $Q(a-v)=\tfrac12\gamma\sigma^2T\left((q-Q)^2-q^2\right)=\gamma\sigma^2T\left(\tfrac{Q^2}{2}-qQ\right)$. The bid follows with $-Q$.
\end{proof}

A dealer who is long shades both quotes down: she wants to sell and discourages buying from others. The inventory term is the root of the market-making
models of One Quant Book 11 (the Avellaneda--Stoikov reservation price has the same form); here it matters that it produces a spread without any
information at all.

\section{Information: the sequential-trade model}

\begin{definition}[Informed trader, noise trader]\label{def:mx:why-there-is-a-spread:informed}
An \emph{informed trader}\index{informed trader} knows something about the value that the liquidity provider does not and trades on it. A \emph{noise
trader}\index{noise trader} (liquidity trader) trades for reasons unrelated to the value (a fund flow, a hedge, a rebalancing) and carries no information.
\end{definition}

\begin{definition}[Sequential-trade model, Glosten--Milgrom model]\label{def:mx:why-there-is-a-spread:gm}
A \emph{sequential-trade model}\index{sequential-trade model} lets traders arrive one at a time to a competitive, risk-neutral liquidity provider who
quotes before each arrival and learns from each trade. The \emph{Glosten--Milgrom model}\index{Glosten--Milgrom model} is the sequential-trade model in
which each arrival is informed with probability $\mu$ and trades in the direction of the value, or is a \omterm{def:mx:why-there-is-a-spread:informed}{noise trader} who buys or sells with equal
probability, and the provider sets the ask and bid to the expected value conditional on a buy and on a sale.
\end{definition}

\begin{omfigure}
\begin{tikzpicture}[font=\footnotesize,grow=right,level distance=2.6cm,
  level 1/.style={sibling distance=2.6cm},level 2/.style={sibling distance=1.25cm},
  every node/.style={align=center}]
\node[draw,rounded corners] {value\\ $v_H$ or $v_L$}
  child {node[draw,rounded corners,fill=omThm!12] {noise\\ $1-\mu$}
    child {node {sells: $\tfrac12$}}
    child {node {buys: $\tfrac12$}}
    edge from parent node[below,pos=0.55] {}}
  child {node[draw,rounded corners,fill=omDef!12] {informed\\ $\mu$}
    child {node {sells if $v=v_L$}}
    child {node {buys if $v=v_H$}}};
\end{tikzpicture}

\omcaption{One arrival of the \omterm{def:mx:why-there-is-a-spread:gm}{Glosten--Milgrom model}. The provider sees only the direction of the trade; a buy is more likely when the value is high,
by $\mu$, and the ask is the value's expectation given a buy.}\label{fig:mx:why-there-is-a-spread:tree}
\end{omfigure}

\begin{proposition}[Glosten--Milgrom quotes]\label{prop:mx:why-there-is-a-spread:gmquotes}
With $v\in\{v_L,v_H\}$, $\Pr(v=v_H)=\delta$ and informed share $\mu$,
\[
a = v_L+(v_H-v_L)\frac{\delta(1+\mu)}{\delta(1+\mu)+(1-\delta)(1-\mu)},\qquad
b = v_L+(v_H-v_L)\frac{\delta(1-\mu)}{\delta(1-\mu)+(1-\delta)(1+\mu)}.
\]
At $\delta=\tfrac12$ the spread is $a-b=\mu(v_H-v_L)$. After a buy the provider's belief becomes $\delta'=a$ (on the scale $v_L=0$, $v_H=1$), and the
quotes are regret-free: they are what each trade reveals.
\end{proposition}

\begin{proof}
$\Pr(\text{buy}\mid v_H)=\mu+\tfrac{1-\mu}{2}=\tfrac{1+\mu}{2}$ and $\Pr(\text{buy}\mid v_L)=\tfrac{1-\mu}{2}$. Bayes' rule (One Quant Book 2, chapter~30,
for the Bayesian update) gives $\Pr(v_H\mid\text{buy})$, and $a=v_L+(v_H-v_L)\Pr(v_H\mid\text{buy})$ by risk neutrality and zero expected profit. The bid is
symmetric. At $\delta=\tfrac12$, $a=v_L+(v_H-v_L)\tfrac{1+\mu}{2}$ and $b=v_L+(v_H-v_L)\tfrac{1-\mu}{2}$.
\end{proof}

\omcode{code/firm/spreadmodels/firm_spreadmodels.py}{29}{39}{Glosten--Milgrom quotes and the update after a trade.}\label{lst:mx:why-there-is-a-spread:gm}

The spread is entirely adverse selection (One Quant Book 1, chapter~1): the provider breaks even across traders, losing to the informed exactly what she
earns from noise. It narrows as trades reveal the value (\cref{fig:mx:why-there-is-a-spread:gm}): with 30\% informed arrivals the median path needs 19
trades for the spread to fall to a tenth of its start, with 10\% it needs 148. Informed trading does two things at once: it widens the spread and makes
prices learn faster. With a processing cost $c$ per trade added to each quote, the adverse-selection share of the spread is $\mu\Delta v/(\mu\Delta
v+2c)$: with $\Delta v=1$ and $c=0.05$, half the spread at $\mu=0.1$ and three quarters at $\mu=0.3$.

\begin{omfigure}
\begin{tikzpicture}
\begin{axis}[width=9cm,height=5.2cm,xlabel={\small arrivals},ylabel={\small mean spread ($v_H-v_L=1$)},tick label style={font=\footnotesize},
  xmin=0,xmax=150,ymin=0,ymax=0.52,grid=major,grid style={gray!25},legend columns=3,legend style={font=\footnotesize,at={(0.5,-0.3)},anchor=north,
  draw=none,fill=none,/tikz/every even column/.append style={column sep=8pt}},table/col sep=comma]
\addplot[omDef,very thick] table[x=n,y=mu10]{figdata/microstructure/04-why-there-is-a-spread/gm.csv};
\addplot[omProp,very thick] table[x=n,y=mu30]{figdata/microstructure/04-why-there-is-a-spread/gm.csv};
\addplot[omThm,very thick] table[x=n,y=mu50]{figdata/microstructure/04-why-there-is-a-spread/gm.csv};
\legend{$\mu=0.1$,$\mu=0.3$,$\mu=0.5$}
\end{axis}
\end{tikzpicture}

\omcaption{The Glosten--Milgrom spread along a sequence of arrivals, averaged over 400 paths (half with a high value): it starts at $\mu$ and falls as
trades reveal the value, faster when more arrivals are informed. Data: \texttt{mx\_spread.gm\_learning}.}\label{fig:mx:why-there-is-a-spread:gm}
\end{omfigure}

\section{Information: the strategic trader}

\begin{definition}[Kyle model]\label{def:mx:why-there-is-a-spread:kyle}
The \emph{Kyle model}\index{Kyle model} has one \omterm{def:mx:why-there-is-a-spread:informed}{informed trader} who knows the value $v\sim\mathcal N(p_0,\Sigma_0)$ and chooses a quantity $x$, \omterm{def:mx:why-there-is-a-spread:informed}{noise
traders} who submit $u\sim\mathcal N(0,\sigma_u^2)$ independent of $v$, and competitive risk-neutral market makers who see only the total order flow
$y=x+u$ and set the price $p=\E[v\mid y]$.
\end{definition}

The \omterm{def:mx:why-there-is-a-spread:informed}{informed trader} is strategic: he knows his trades move the price and hides in the noise. The linear equilibrium is short to derive.

\begin{proposition}[Kyle's one-period equilibrium]\label{prop:mx:why-there-is-a-spread:kyle}
In the unique linear equilibrium $x=\beta(v-p_0)$ and $p=p_0+\lambda y$ with
\[
\beta=\frac{\sigma_u}{\sqrt{\Sigma_0}},\qquad \lambda=\frac{\sqrt{\Sigma_0}}{2\sigma_u};
\]
the price reveals half of the private information ($\mathrm{Var}(v\mid p)=\Sigma_0/2$) and the \omterm{def:mx:why-there-is-a-spread:informed}{informed trader} earns $\tfrac12\sqrt{\Sigma_0}\,\sigma_u$ in
expectation, paid by the \omterm{def:mx:why-there-is-a-spread:informed}{noise traders}.
\end{proposition}

\begin{proof}
Given $p=p_0+\lambda y$, the \omterm{def:mx:why-there-is-a-spread:informed}{informed trader} maximises $\E[x(v-p)\mid v]=x(v-p_0)-\lambda x^2$, so $x=(v-p_0)/(2\lambda)$ and $\beta=1/(2\lambda)$. Given
$x=\beta(v-p_0)$, the market maker's projection gives $\lambda=\mathrm{Cov}(v,y)/\mathrm{Var}(y)=\beta\Sigma_0/(\beta^2\Sigma_0+\sigma_u^2)$. Substituting
$\beta=1/(2\lambda)$ gives $\beta^2\Sigma_0=\sigma_u^2$, hence both formulas. Then $\mathrm{Var}(v\mid y)=\Sigma_0-\lambda^2\mathrm{Var}(y)=\Sigma_0-\Sigma_0/2$,
and the expected profit is $\beta\Sigma_0-\lambda\beta^2\Sigma_0=\sqrt{\Sigma_0}\sigma_u/2$.
\end{proof}

The slope $\lambda$ is Kyle's lambda (One Quant Book 7, chapter~9): the price impact of a unit of order flow, the inverse of the market's depth. More noise
makes the market deeper and the insider richer, and leaves the price exactly as informative. On 100\,000 simulated draws with $\sqrt{\Sigma_0}=1$ and
$\sigma_u=2$, the regression of $p$ on $y$ returns 0.250, the residual variance of $v$ is 0.498 and the insider's average profit 0.991, against the
proposition's 0.25, 0.5 and 1. In Kyle's continuous-time version the insider trades gradually, $\lambda$ is constant through time and all of his
information is in the price by the end (\admitted; derived in Kyle's paper).

\section{Measuring informed trading}

\begin{definition}[Probability of informed trading]\label{def:mx:why-there-is-a-spread:pin}
The \emph{probability of informed trading}\index{probability of informed trading} (PIN) of Easley, Kiefer, O'Hara and Paperman is, in their model of daily
order flow, the share of trades that come from \omterm{def:mx:why-there-is-a-spread:informed}{informed traders}: $\mathrm{PIN}=\alpha\mu/(\alpha\mu+\varepsilon_b+\varepsilon_s)$, where an information event
occurs on a day with probability $\alpha$, is bad news with probability $\delta$, brings informed orders at rate $\mu$, and uninformed buys and sells arrive
at rates $\varepsilon_b$ and $\varepsilon_s$.
\end{definition}

Each day's buys and sells are Poisson with rates set by the day's hidden state, and the parameters are fitted by maximum likelihood on daily counts (the
likelihood mixes three Poisson pairs; its logarithm must be computed in a numerically stable way).

\omcode{code/firm/spreadmodels/firm_spreadmodels.py}{106}{113}{The PIN log-likelihood, stable through log-sum-exp.}\label{lst:mx:why-there-is-a-spread:pin}

The model reads any common swing in buying and selling as information when the swing is not one-sided in its data, and Duarte and Young (2009) found that
PIN's association with returns came from its illiquidity component rather than its information component.
\Cref{fig:mx:why-there-is-a-spread:pin} shows why on simulated data. Days with no information at all, whose uninformed rates move together from day to day
(a lognormal factor on both), are fitted with a PIN of 0.085 when the factor's standard deviation is 0.25, 0.153 at 0.5 and 0.207 at 0.75, against 0.006
without the swings. A market with true PIN 0.107 is estimated at 0.107 without swings and 0.195 with the largest.

\begin{omfigure}
\begin{tikzpicture}
\begin{axis}[width=9cm,height=5.2cm,xlabel={\small standard deviation of the daily activity factor},ylabel={\small fitted PIN},
  tick label style={font=\footnotesize},xtick={0,0.25,0.5,0.75},ymin=0,ymax=0.25,grid=major,grid style={gray!25},
  legend columns=2,legend style={font=\footnotesize,at={(0.5,-0.3)},anchor=north,draw=none,fill=none,
  /tikz/every even column/.append style={column sep=8pt}},table/col sep=comma]
\addplot[omDef,very thick,mark=*,error bars/.cd,y dir=both,y explicit] table[x=sd,y=none,y error=none_sd]{figdata/microstructure/04-why-there-is-a-spread/pin.csv};
\addplot[omThm,very thick,mark=square*,error bars/.cd,y dir=both,y explicit] table[x=sd,y=info,y error=info_sd]{figdata/microstructure/04-why-there-is-a-spread/pin.csv};
\addplot[omThm,dashed,domain=0:0.75] {0.1071};
\legend{no information (true PIN 0),information (true PIN 0.107)}
\end{axis}
\end{tikzpicture}

\omcaption{Maximum-likelihood PIN on simulated samples of 250 days (mean and standard deviation over 20 samples), against the size of a common daily
factor on the uninformed rates; the dashed line is the true PIN of the informed market. Data: \texttt{mx\_spread.pin\_bias}.}\label{fig:mx:why-there-is-a-spread:pin}
\end{omfigure}

\section{Tutorial: three spreads in code}\label{tut:mx:why-there-is-a-spread:tut}

\begin{tutorial}
\textbf{Goal.} Compute the three spreads, check Kyle's equilibrium by simulation and see PIN mistake activity for information.
\textbf{End state:} \cref{fig:mx:why-there-is-a-spread:gm,fig:mx:why-there-is-a-spread:pin} and the numbers of sections 1 and 4.
\begin{tutsteps}
\item \textbf{Bounce.} \texttt{mx\_spread.roll\_on\_tape()}: Roll's estimate, the quoted spread and the trade-sign autocorrelation of the simulated hour.
\item \textbf{Inventory.} \texttt{inventory\_quotes(v, q, size, gamma, sigma, horizon)} for inventories from $-5$ to $5$: the quotes move against the
position.
\item \textbf{Learning.} \texttt{gm\_learning()} and \texttt{gm\_with\_cost(mu, dv, c)}.
\item \textbf{Kyle.} \texttt{kyle\_check()} simulates the equilibrium and recovers $\lambda$.
\item \textbf{PIN.} \texttt{pin\_bias(sd, alpha)} over the grid; draw with \texttt{fig\_spread.py}.
\end{tutsteps}
\textbf{What to change next.} Let the \omterm{def:mx:why-there-is-a-spread:informed}{informed trader} in the \omterm{def:mx:why-there-is-a-spread:gm}{Glosten--Milgrom model} trade only when the quote is far enough from the value (a profit
threshold) and watch the spread; fit PIN on \texttt{firm.tape}'s days of buys and sells, where the informed flag is known.
\end{tutorial}

\section{Build: the spread models}\label{bld:mx:why-there-is-a-spread:spreadmodels}

\begin{build}
\bfield{Purpose} The reference implementations of the spread theories, used by chapter~5's decompositions, by the agent-based market of chapter~27
(an informed agent and a Glosten--Milgrom market maker) and by Book~11's market-making tutorials as benchmarks.
\bfield{Interface} \texttt{gm\_quotes}, \texttt{gm\_update}, \texttt{gm\_path}; \texttt{kyle\_one\_period}, \texttt{kyle\_simulate};
\texttt{inventory\_quotes}; \texttt{roll\_spread}; \texttt{pin}, \texttt{pin\_loglik}, \texttt{pin\_fit}, \texttt{simulate\_days}.
\bfield{Rules} Closed forms where they exist; the PIN likelihood in log-sum-exp form with Poisson log-probabilities from \texttt{gammaln}; fits from several
starting points; everything seeded.
\bfield{Acceptance tests} \texttt{code/firm/spreadmodels/tests/}: the Glosten--Milgrom spread equal to $\mu$ at $\delta=\tfrac12$ and the ask equal to the
posterior after a buy; Kyle's $\lambda$, $\beta$, posterior variance and profit, and their simulated values; inventory quotes by hand; Roll's estimator on
a bouncing random walk; PIN recovered within 0.02 on 300 simulated days.
\bfield{Stretch} Kyle's multi-period recursion; the Easley--O'Hara model with time-varying arrival rates; the adjusted PIN of Duarte and Young.
\end{build}

\begin{omsources}
\item R.~Roll, ``A simple implicit measure of the effective bid-ask spread in an efficient market'', \emph{Journal of Finance} 39(4), 1984.
\item M.~B.~Garman, ``Market microstructure'', \emph{Journal of Financial Economics} 3(3), 1976.
\item H.~R.~Stoll, ``The supply of dealer services in securities markets'', \emph{Journal of Finance} 33(4), 1978.
\item T.~Ho and H.~R.~Stoll, ``Optimal dealer pricing under transactions and return uncertainty'', \emph{Journal of Financial Economics} 9(1), 1981.
\item L.~R.~Glosten and P.~R.~Milgrom, ``Bid, ask and transaction prices in a specialist market with heterogeneously informed traders'', \emph{Journal of
Financial Economics} 14(1), 1985.
\item A.~S.~Kyle, ``Continuous auctions and insider trading'', \emph{Econometrica} 53(6), 1985.
\item D.~Easley, N.~M.~Kiefer, M.~O'Hara and J.~B.~Paperman, ``Liquidity, information, and infrequently traded stocks'', \emph{Journal of Finance} 51(4), 1996.
\item J.~Duarte and L.~Young, ``Why is PIN priced?'', \emph{Journal of Financial Economics} 91(2), 2009.
\end{omsources}

\section{Exercises}

\begin{exercise}[$\star$]\label{exo:mx:why-there-is-a-spread:1}
Successive transaction-price changes of a stock have an autocovariance of $-0.0004$ (dollars squared). What spread does Roll's estimator give?
\end{exercise}

\begin{exercise}[$\star$]\label{exo:mx:why-there-is-a-spread:2}
In the \omterm{def:mx:why-there-is-a-spread:gm}{Glosten--Milgrom model} the value is 99 or 101 with equal probability and 20\% of arrivals are informed. What are the bid and the ask?
\end{exercise}

\begin{exercise}[$\star$]\label{exo:mx:why-there-is-a-spread:3}
In Kyle's model $\sqrt{\Sigma_0}=2$ and $\sigma_u=8$. Give $\lambda$, $\beta$, the insider's expected profit and the posterior variance of the value.
\end{exercise}

\begin{exercise}[$\star\star$]\label{exo:mx:why-there-is-a-spread:4}
A dealer marks her position to 50, is short 3 units, and has $\gamma=0.2$, $\sigma=1$ and $T=0.5$. What are her bid and ask for one unit? Explain why
both are above 50.
\end{exercise}

\begin{exercise}[$\star\star$]\label{exo:mx:why-there-is-a-spread:5}
In exercise~2, a buy arrives. What is the provider's new probability of the high value, and what are the new quotes?
\end{exercise}

\begin{exercise}[$\star\star$]\label{exo:mx:why-there-is-a-spread:6}
Why does doubling the \omterm{def:mx:why-there-is-a-spread:informed}{noise traders}' volume in Kyle's model not make the price less informative?
\end{exercise}

\begin{exercise}[$\star\star\star$]\label{exo:mx:why-there-is-a-spread:7}
\emph{Coding.} Fit PIN on 250 simulated days with $\alpha=0.4$ and a daily activity factor of standard deviation 0.5, then fit it again after dividing each
day's buys and sells by the day's total and multiplying by the sample's mean total. Does normalising remove the bias? Why only partly?
\end{exercise}

\begin{exercise}[$\star\star\star$]\label{exo:mx:why-there-is-a-spread:8}
\emph{Find the flaw.} ``Stock A has a PIN of 0.25 and stock B of 0.12, so A's market makers face twice as much informed trading.''
\end{exercise}

\section{Problem: The Specialist's Two Cents}

\begin{problem}[{Weekend problem --- the specialist's two cents}]\label{pb:mx:why-there-is-a-spread:1}
A specialist quotes a two-cent spread and wants to know what it pays for: processing, inventory or information; and a research desk wants to know whether
PIN can tell.

\medskip\noindent\textbf{Part I --- Processing and inventory.}
\begin{enumerate}
\item What does Roll's estimator assume, and what does it give on the simulated hour?
\item Why does it fall short of the quoted spread there?
\item State the mean--variance dealer's quotes and the spread they imply.
\item What happens to the quotes when the dealer is long, and why?
\end{enumerate}

\medskip\noindent\textbf{Part II --- The \omterm{def:mx:why-there-is-a-spread:gm}{sequential-trade model}.}
\begin{enumerate}[resume]
\item Derive the Glosten--Milgrom ask at $\delta=\tfrac12$.
\item What is the spread with 30\% informed arrivals and a value that is 0 or 1?
\item How many trades does the median path need for the spread to fall to a tenth, at 10\%, 30\% and 50\% informed?
\item With a processing cost of 0.05 per trade, what share of the spread is adverse selection at $\mu=0.1$ and $\mu=0.3$?
\end{enumerate}

\medskip\noindent\textbf{Part III --- The strategic trader.}
\begin{enumerate}[resume]
\item Derive Kyle's $\lambda$ and $\beta$.
\item What share of the private information does the price reveal?
\item What do 100\,000 simulated draws give for $\lambda$, the posterior variance and the insider's profit?
\item Who pays the insider's profit?
\end{enumerate}

\medskip\noindent\textbf{Part IV --- PIN.}
\begin{enumerate}[resume]
\item Write PIN in terms of the model's parameters and compute it for $\alpha=0.4$, $\mu=60$, $\varepsilon_b=\varepsilon_s=100$.
\item What PIN is fitted when there is no information and no activity swing?
\item What is fitted with no information and daily activity swings of standard deviation 0.25, 0.5 and 0.75?
\item State the \emph{named result}: the adverse-selection share of the spread against the informed share, and the bias of PIN under common activity swings.
\item Why does the model read common swings as information?
\item What did Duarte and Young conclude about PIN and returns?
\item How would you separate illiquidity from information in daily order flow?
\item In one sentence: which of the three costs does a two-cent spread in a liquid stock mostly pay for?
\end{enumerate}
\end{problem}

\section{Interview questions}

\begin{interviewq}[$\star$ \iqroles{trader, researcher}]\label{iq:mx:why-there-is-a-spread:1}
Name the three components of the bid--ask spread and one piece of evidence for each.
\end{interviewq}

\begin{interviewq}[$\star\star$ \iqroles{researcher}]\label{iq:mx:why-there-is-a-spread:2}
Derive the Glosten--Milgrom ask price when the value is 0 or 1 with equal probability and a fraction $\mu$ of traders is informed.
\end{interviewq}

\begin{interviewq}[$\star\star$ \iqroles{researcher}]\label{iq:mx:why-there-is-a-spread:3}
In Kyle's model, what happens to the price impact and to the insider's profit when noise trading doubles?
\end{interviewq}

\begin{interviewq}[$\star\star$ \iqroles{trader}]\label{iq:mx:why-there-is-a-spread:4}
You make a market and you are long 10\,000 shares. How do you move your quotes, and why?
\end{interviewq}

\begin{interviewq}[$\star\star$ \iqroles{researcher}]\label{iq:mx:why-there-is-a-spread:5}
Roll's estimator gives an imaginary spread for a stock. What happened, and what do you do?
\end{interviewq}

\begin{interviewq}[$\star\star\star$ \iqroles{researcher}]\label{iq:mx:why-there-is-a-spread:6}
How would you estimate the share of a stock's spread that compensates adverse selection?
\end{interviewq}
